% Slide 08: View-based 3D retrieval
\begin{frame}[c]{Rendered views make 3D models searchable}
\centering
\begin{tikzpicture}[node distance=0.52cm, every node/.style={font=\scriptsize}]
  \node[redbox, minimum width=1.7cm, minimum height=0.8cm] (mesh) {3D mesh};
  \node[bluebox, right=of mesh, minimum width=2.0cm, minimum height=0.8cm] (views) {multiple\\camera views};
  \node[mintbox, right=of views, minimum width=2.0cm, minimum height=0.8cm] (renders) {2D renders};
  \node[bluebox, right=of renders, minimum width=2.0cm, minimum height=0.8cm] (clip) {CLIP image\\encoder};
  \node[greybox, right=of clip, minimum width=2.0cm, minimum height=0.8cm] (vectors) {view\\descriptors};

  \draw[flow] (mesh) -- (views);
  \draw[flow] (views) -- (renders);
  \draw[flow] (renders) -- (clip);
  \draw[flow] (clip) -- (vectors);

  \node[mintbox, below=0.95cm of vectors, minimum width=2.0cm] (parent) {linked GLB model};
  \draw[flow] (vectors) -- node[right,tinylabel] {linked to} (parent);

\end{tikzpicture}

\vspace{0.45cm}
\begin{beamercolorbox}[sep=0.5em,rounded=true]{palette highlight}
\small
Rather than encoding the mesh itself, the system searches over rendered views.
A matching view points back to the original GLB model, which is then loaded
into the museum.
\end{beamercolorbox}
\end{frame}
